\documentclass[11pt]{article}
\usepackage[letterpaper,margin=.94in]{geometry}
\usepackage{amsmath,amssymb,amsthm,booktabs,tabularx,microtype}
\usepackage[T1]{fontenc}
\usepackage{lmodern}
\usepackage[colorlinks=true,linkcolor=teal,urlcolor=teal,citecolor=teal]{hyperref}
\usepackage{xcolor}
\newtheorem{theorem}{Theorem}
\newtheorem{lemma}[theorem]{Lemma}
\newcommand{\F}{\mathbb F}
\newcommand{\Z}{\mathbb Z}
\newcommand{\Q}{\mathbb Q}
\newcommand{\ord}{v_p}
\newcommand{\tri}[1]{\tfrac{#1(#1+1)}2}
\newcommand{\repo}[2]{\href{https://github.com/VladimirReshetnikov/ProveIt/blob/6d284fcf75bc3a91bb2e79ba7636f94724bff62c/Analysis/GeneralizedRiemannHypothesis/#1}{#2}}
\title{Sharp root counts for sparse polynomials\\in the open $p$-adic exponential disk}
\author{Research supplement to the ProveIt GRH project}
\date{27 November 2026}
\begin{document}
\maketitle
\begin{abstract}
An integer exponent has the same residue after any positive iterate of Frobenius. We use this observation and a filtration of the exponent algebra to count the roots of a polynomial in a multiplicative disk of relative depth strictly greater than $1/(p-1)$. For any fixed support of size $t\le p$, its optimal bound is determined by the multiplicities of its exponents modulo $p$. Maximizing over supports gives $(h-1)p$, with an extra root exactly when $t+1$ is triangular, where $h$ is the largest integer with $h(h+1)/2\le t$. Explicit products attain the bound with simple roots. The uniform formula continues through $t=p+1$. At $t=p+2$ the maximum jumps to $p^2$. All results are unconditional.
\end{abstract}

\section{The extremal problem}
Write $\ord(p)=1$ on $\overline{\Q}_p$ and put
\[
 D(A)=\{x:\ord(x/A-1)>1/(p-1)\},\qquad A\in\overline{\Q}_p^\times.
\]
Roots are geometric and counted with multiplicity. This convention permits arbitrary finite ramification in both coefficients and roots. With integer exponents, multiplying by a monomial preserves this count. Scaling $x$ by $A$ reduces the extremal problem to $A=1$.

\begin{theorem}[Sharp bound in terms of the number of terms]\label{thm:all}
Let $2\le t\le p$ and write $t=h(h+1)/2+e$ with $0\le e\le h$. Every nonzero polynomial with at most $t$ terms has at most
\[
 B_p(t)=(h-1)p+\mathbf 1_{e=h}
\]
roots in $D(A)$. An integer polynomial of degree $B_p(t)$ achieves equality with distinct roots about $1$. For $p>2$ the same extremum holds on every closed disk of relative depth $r$ in
\[
 1/(p-1)<r\le 1/(p-2).
\]
For $p=t=2$ the maximum is one at every greater depth.
\end{theorem}

For example, subject to $t\le p$, the successive maxima for $t=6,7,8,9,10$ are $2p,2p,2p,2p+1,3p$. With $t=p$ and $p\to\infty$ the answer is $\sqrt{2}\,p^{3/2}+O(p)$. A closed disk at the boundary itself can contain additional roots. The strict inequality in $D(A)$ is essential.

\section{The complete invariant for a fixed support}
Let $S$ be a set of $t$ distinct nonnegative exponents. For each residue $a\in\F_p$, let $s_a=\#\{n\in S:n\equiv a\pmod p\}$. Form a diagram with row lengths
\[
 m_b=\#\{a:s_a>b\},\qquad b=0,1,\ldots.
\]
Its boxes are $(i,b)$ with $0\le i<m_b$. Define the length of the degree-$d$ diagonal by $H_d=\#\{(i,b):i+b=d\}$. The diagram has $t$ boxes.
\begin{theorem}[Fixed-support formula]\label{thm:fixed}
For $1\le |S|=t\le p$ the maximum number of roots in $D(A)$ over nonzero polynomials supported on $S$ is
\[
 C_p(S)=\max_{d:H_d>0}\bigl(d+(p-1)(H_d-1)\bigr).
\]
The upper bound holds with arbitrary algebraic $p$-adic coefficients. At $A=1$ one may choose integer coefficients, every exponent in $S$ present, and all the roots counted in $D(1)$ distinct.
\end{theorem}
For instance, the supports $\{0,1,\ldots,5\}$ and $\{0,p,2p,1,p+1,2\}$ both have six elements when $p\ge7$. The first has row lengths $(6)$ and maximum five. The second has $(3,2,1)$ and diagonal lengths $1,2,3$, giving maximum $2p$. The actual representatives in each residue class can be arbitrarily large.

\subsection*{The analytic reduction}
Factor out $A$ by the preceding scaling and fix a rational $r>1/(p-1)$. Choose $z$ with $\ord(z)=r$. On $\ord(Y)\ge0$, the substitution $X=\exp(zY)$ identifies the closed disk with $\ord(X-1)\ge r$. Its inverse is $\log(X)/z$ and its derivative never vanishes. For $P=\sum_{n\in S} c_n X^n$ the resulting series has coefficients
\[
 M_k z^k/k!,\qquad M_k=\sum_{n\in S} c_n n^k .
\]
Put $A_0=\Z_{(p)}[T]/\prod_{n\in S}(T-n)$ and define $L(T^k)=M_k$. The integral monic relation shows that every $M_k$ is an integral linear combination of the first $t$. The Vandermonde matrix on the distinct $n$ is nonsingular over $\Q$. After division by a nonzero moment of least valuation among the first $t$, $L$ takes integral values on $A_0$ and some early moment is a unit.

Set $s(k)$ equal to the sum of the base-$p$ digits of $k$ and $\delta=r-1/(p-1)>0$. Legendre's formula gives
\[
 w(k)=rk-\ord(k!)=\delta k+s(k)/(p-1).
\]
The coefficients tend to zero. Weierstrass preparation in the one-variable Tate algebra identifies the number of geometric roots on the closed disk with the largest index minimizing $\ord(M_k)+w(k)$, taking $\ord(0)=+\infty$. The minimum is at most $w(t-1)$ since some early moment is a unit and $w(i)=ri$ for $i<t$. Thus at a minimizing $k\ge t$ we have $s(k)\le t-2$ and $\ord(M_k)<1$. Indeed, digit sum at least $t-1$ would make $w(k)>w(t-1)$, and for $k\ge t$ the difference $w(t-1)-w(k)$ is strictly less than one.

\subsection*{A basis ordered by digit sum}
In $R=A_0/pA_0$ put $U=T^p$ and $V=T-U$. The Chinese remainder theorem gives factors $\F_p[T]/(T-a)^{s_a}$ for $s_a>0$. In each factor $U=a$ and $V=T-a$, since $s_a\le p$. All powers $T^{p^j}$ for $j\ge1$ therefore equal $U$. If $k$ has digit sum $d$ and its digits above the units place sum to $j$, then
\begin{equation}\label{eq:compress}
 T^k=T^{d-j} U^j \text{ in }R,\qquad d+(p-1)j\le k .
\end{equation}
Filter $R$ by polynomial total degree in $U,V$. The monomials $U^iV^b$ indexed by diagram boxes are a basis adapted to this filtration. In fact,
\[
 V^b\prod_{a:s_a>b}(U-a)=0 .
\]
The monic relation for each $b$ reduces every monomial in its row to diagram boxes of no greater degree. There are $t$ boxes, and the spanning family therefore is a basis since $\dim R=t$ and $T=U+V$ generates $R$. Each reduction of a monomial outside the diagram strictly lowers its degree.

For any nonempty diagonal of degree $d$, its highest possible row index is at most $d$, and $d\le t-1<p$. In degree $d$ of the associated graded algebra consider the $H_d$ columns
\[
 T^{d-j} U^j,\qquad 0\le j<H_d .
\]
They form a basis. In the original box basis their coefficient in row $b$ is $\binom{d-j}{b}$. With $(x)_j=x(x-1)\cdots(x-j+1)$ and $(x)_0=1$ one has
\[
 \binom{d-j}{b}=\binom db\,(d-b)_j/(d)_j .
\]
All row and column scale factors are units modulo $p$. The remaining matrix evaluates monic polynomials of successive degrees at the distinct residues $d-b$ and its Vandermonde determinant is a unit. This proves the assertion, including the constant diagonal.

Consequently every class of filtered degree $d$ is spanned by the displayed columns in degrees $d'\le d$. Their actual power indices $\ell=d'+(p-1)j$ have digit sum $d'$, since $0\le j\le d'<p$. They are distinct by uniqueness of base-$p$ digits. These $t$ indices lie between zero and $C_p(S)$, so $C_p(S)\ge t-1$.

Suppose a minimizing index $k$ exceeded $C_p(S)$. Then $k\ge t$. By \eqref{eq:compress} we can write $T^k$ modulo $p$ as a linear combination of the basis powers with $\ell\le C_p(S)$ and $s(\ell)\le s(k)$. All have $w(\ell)<w(k)$. Lift the identity to $A_0$; the error is in $p A_0$, on which $L$ has valuation at least one. Since $\ord(M_k)<1$, some participating moment has valuation at most $\ord(M_k)$. Its coefficient has smaller valuation, a contradiction. This argument works unchanged in ramified coefficient fields. It bounds the roots on every rational closed disk strictly inside $D(1)$, hence also on $D(1)$ since a polynomial has only finitely many roots.

\subsection*{Sharpness for each prescribed support}
Write $C=C_p(S)$. If $C=t-1$, Vandermonde supplies a nonzero rational functional with $M_0=\cdots=M_{t-2}=0$ and $M_{t-1}=1$. Its corresponding polynomial has a zero of multiplicity $t-1$ at $1$. The monomial case $t=1$ follows immediately. Assume henceforth $C\ge t$ and choose $d,j$ attaining the maximum with $j=H_d-1$. Then $d\le p-2$. Indeed a box of degree $p-1$ requires at least $p$ boxes and, when $t\le p$, can only occur in a single row or a single column: its southwest rectangle has $(i+1)(b+1)=p+ib$ boxes. In these exceptional diagrams $C=p-1$.

Choose an $\F_p$-linear functional on $R$ annihilating all preceding-degree columns and the earlier columns of degree $d$, and taking the value $1$ on the column with index $C$. Lift its values on $1,T,\ldots,T^{t-1}$ to integers to get $L$ on $A_0$. Rational Vandermonde inversion gives rational $c_n$ with $M_k=L(T^k)$ for all $k$. Choose a rational $r$ sufficiently close to the critical value from above so that
\[
 w(C)<\min\{1,(d+1)/(p-1)\}.
\]
If $w(k)<w(C)$ then $s(k)\le d$. For $s(k)<d$ its residue class lies in the annihilated filtration. For $s(k)=d$, one has $k<C$, so \eqref{eq:compress} gives a column with a strictly smaller $j$ and again $M_k\equiv0\pmod p$. Those coefficients have valuation at least $1$. All other coefficients have valuation at least $w(C)$, and the $C$-th has valuation exactly $w(C)$. Preparation gives at least $C$ geometric zeros on this closed disk. The upper bound gives equality.

In either case a sufficiently small rational perturbation of each permitted coefficient preserves the minimum valuation and its last index on the chosen disk: all moment perturbations have valuation at least the minimum of the coefficient perturbations, and $w(k)\ge0$. We may first ensure every coefficient is nonzero. After factoring out the smallest power of $X$, vary the constant slightly to avoid the finitely many critical values of the resulting nonconstant polynomial. This makes all nonzero roots simple. Clearing denominators completes the proof of Theorem~\ref{thm:fixed}.

\section{Optimization over supports}
Suppose a diagram has $s\ge1$ boxes on diagonal $d$, at rows $b_1<\cdots<b_s$. Set $g_0=b_1$, $g_s=d-b_s$, and $g_i=b_{i+1}-b_i-1$ for $1\le i<s$. Each is nonnegative and their sum is $d+1-s$. The diagram contains the union of the southwest rectangles below these boxes. Its area is
\[
 \begin{aligned}
 \sum_{i=1}^{s}(1+g_{i-1})\Big(s-i+1+\sum_{j=i}^{s}g_j\Big)
 &=\frac{s(s+1)}2+s(d+1-s)+\sum_{0\le i<j\le s}g_i g_j .
 \end{aligned}
\]
Indeed the successive row blocks have the heights and widths in the two factors. Thus $s\le h$. If $s=h$, we obtain $d\le h-1+\lfloor e/h\rfloor$. If $s<h$, use $d\le t-1\le p-1$ to obtain $d+(p-1)(s-1)\le(p-1)(h-1)$. Together these imply $C_p(S)\le B_p(t)$. For fewer than $t$ exponents one may pad the support by unused columns.

For matching lower bounds take $p>2$ and $m=h-1$. For $1\le i\le m$ put
\[
 G_i(X)=X^p-pX+(p-1)+i p^3,\qquad
 F(X)=(X-1)^{\mathbf1_{e=h}}\prod_{i=1}^m G_i(X).
\]
The exponents in the product are among $a+pb$ with $a,b\ge0$, $a+b\le m$. There are at most $(m+1)(m+2)/2$ of them. Multiplication by $X-1$ adds at most $m+1=h$ further exponents (the pairs of total degree $m+1$ with first coordinate positive). Hence $F$ meets the term budget and its degree is exactly $B_p(t)$.

For $Y=X-1$ the constant coefficient of $G_i(1+Y)$ has valuation at least three, its linear coefficient vanishes, and its coefficients in degrees $2,\ldots,p-1$ have valuation one. If $\ord(Y)<1/(p-2)$ the monic term of degree $p$ has strictly smaller valuation than every other term: use $3\ge p/(p-2)$ for the constant. Thus all $p$ roots have relative depth at least $1/(p-2)$. Two factors differ by a nonzero constant. A multiple root of $G_i$ would satisfy $X^{p-1}=1$ and $X=1+i p^3/(p-1)$, which are incompatible. Also $G_i(1)\ne0$. This proves all the sharpness statements, with $X-1$ covering $t=2$, including $p=2$.

\section{The first jump beyond the residue characteristic}
\begin{theorem}[The $p+2$ threshold]\label{thm:jump}
For every prime $p$, the sharp open-disk bound for at most $p+1$ terms is $B_p(p+1)$ as given by the same triangular formula. For at most $p+2$ terms the exact maximum is $p^2$. In both cases integer polynomials achieve the bound with distinct roots.
\end{theorem}
At $p=11$ this gives 33 roots with twelve permitted terms and 121 with thirteen. As $p$ grows, the extra term increases the extremum by a factor asymptotic to $\sqrt{p/2}$.

\begin{proof}
As before we work about $1$. For a support of size $t\ge2$ whose exponents all have the same residue, factor out its smallest monomial and divide all exponent differences by their common power $p^k$, so that the polynomial takes the form $X^a Q(X^{p^k})$. Exponentiation by $p^k$ maps $D(1)$ injectively into itself, via $\log$, and preserves multiplicities since its derivative is nonzero. It therefore suffices for upper bounds to consider supports occupying at least two residues. Polynomials with fewer terms are already covered by the earlier cases. Work with exact $t=p+1$ and then $t=p+2$.

For $p\ge3$, all integers $i<t\le p+2$ have $s(i)\le p-1$. Consequently any minimizing index $k\ge t$ has $s(k)\le p-2$ and $\ord(M_k)<1$, by comparison with the unit among the first $t$ normalized moments. In detail, $w(i)-w(k)\le\delta(t-1-k)+(p-1-s(k))/(p-1)$. For $p=2,t=3$ the initial digit sums are at most one, so all $k\ge3$ are excluded immediately. For $p=2,t=4$ the initial sums are at most two; the same comparison gives $s(k)\le1$ and $\ord(M_k)<1$ for any minimizing $k\ge4$.

If $t=p+1$, every residue has multiplicity at most $p$. When $p\ge3$ the preceding diagram proof applies to all relevant filtered degrees $d\le p-2$, so all relevant power indices are at most $B_p(t)$ by the southwest area bound. The initial indices are at most $t-1=p\le B_p(t)$. The product construction gives matching witnesses at odd primes, since its term-budget calculation required only $p>2$. For $p=2$, the polynomial $(X-1)(X-5)$ gives two distinct roots of depth greater than one.

Now let $t=p+2$. If every residue has multiplicity at most $p$, \eqref{eq:compress} applies. For each potentially minimizing $k\ge t$ its compressed index is at most $p\,s(k)<p^2$; this also holds in the binary case. When $k>p^2$ the weight decreases strictly under compression, and lifting the residue identity rules out $k$. The only other multiplicity profile is $(p+1,1)$. Subtract a common exponent from the support, or add a common nonnegative exponent, to place the large class at zero. This monomial shift preserves roots on the disk. In the residue algebra $\F_p[T]/T^{p+1}\times\F_p$ the second factor evaluates $T$ at some $a\ne0$. For every $k>p^2$ one has
\[
 T^k=a^{k-p^2}T^{p^2},
\]
because both powers vanish on the first factor. Since $s(k)\ge1=s(p^2)$, the index $p^2$ has strictly smaller weight; lifting again excludes $k$. This proves the upper bound.

For sharpness use
\[
 P(X)=(X^p-1)^p-p^p\bigl(X^p-1-p(X-1)\bigr).
\]
Its support is contained in $\{0,p,2p,\ldots,p^2\}\cup\{1\}$. Choose $\pi$ with $\pi^{p-1}=p$ and put $X=1+\pi Z$. The polynomial $H(Z)=((1+\pi Z)^p-1)/\pi^p$ has integral coefficients and reduces to $Z+Z^p$ at the maximal ideal. Hence
\[
 \pi^{-p^2}P(1+\pi Z)=H(Z)^p-H(Z)+Z\equiv Z^{p^2}.
\]
This is a monic integral polynomial, so all its roots are integral; their reductions are zero. Thus all $p^2$ roots of $P$ lie strictly inside $D(1)$. Adding a sufficiently $p$-divisible rational integer to its constant term preserves the last congruence and can avoid the finitely many critical values, giving $p^2$ distinct roots. 
\end{proof}

\clearpage
\section{Relation to the source and further questions}
The repository's \repo{Article/sections/smd.tex}{local root-count lemma} derives the factorial-weight bound from the moment recurrence. The \repo{Research/padic-fewnomial-root-count/report.tex}{two-to-five-term report} adds digit compression and proves the uniform bound $p(t-2)$ for $3\le t\le p$; the \repo{Research/padic-six-term/report.tex}{six-term report} treats all closed-disk depths at $t=6$, $p\ge7$. Their near-boundary values agree with Theorem~\ref{thm:all}. The diagonal basis above replaces the small-case residue profiles and yields both the all-$t$ formula and the full invariant of a prescribed support. The independence from the ambient GRH input follows directly from the proof. These theorems have not been kernel-checked here.

Van der Poorten's \href{https://doi.org/10.1017/S144678870001332X}{1976 Hermite interpolation theorem} bounds zeros of more general exponential polynomials as a function of a frequency valuation margin. Van der Poorten and Rumely treat the entire convergence disk in \href{https://doi.org/10.1112/jlms/s2-36.1.1}{\emph{Zeros of $p$-adic exponential polynomials II}} (1987). Lenstra's \href{https://doi.org/10.1515/9783110285581.277}{local fewnomial bounds} also include results counting zeros in a specified finite field extension. The fixed-support formula concerns geometric roots of integer-frequency sums. The comparison establishes the extension of the cited repository results; publication priority is unresolved.

\subsection*{Independent finite checks}
The accompanying Python script derives the root counts of the explicit witnesses directly from the coefficients of $P(1+Y)$. It minimizes $\ord([Y^k]P(1+Y))+rk$ over $k$. At the critical radius $r=1/(p-1)$ a slight increase in $r$ breaks ties in favor of the smallest minimizing index, giving the open-disk count. Each displayed test for simplicity uses a squarefree reduction at an auxiliary prime.

\begin{center}
\begin{tabular}{r@{\hspace{1.9em}}rrrrrr}
\toprule
$p$ & 3 & 5 & 7 & 11 & 13 & 17\\
\midrule
maximum with $p+1$ terms & 3 & 10 & 14 & 33 & 40 & 68\\
maximum with $p+2$ terms & 9 & 25 & 49 & 121 & 169 & 289\\
\bottomrule
\end{tabular}
\end{center}
The separate residue-algebra check enumerates multiplicity partitions. It computes jets of $T^k$ in each residue factor as $\binom{k}{b}a^{k-b}$, then scans powers by digit sum and compressed index using Gaussian elimination. The first independent columns agree with the claimed diagonal bases in all 1,274 tested assignments. Run \texttt{python3 verify.py} to reproduce the checks recorded in \texttt{verification.txt}. The proofs above apply at every prime.

\subsection*{Further research questions}
Can one classify the changes in $C_p(S)$ on deeper closed disks for arbitrary $t$? For larger $t$ the Frobenius power $T^p$ can retain a nonzero nilpotent part on a residue factor. Which additional residue data determine the sharp bound in that regime?
\end{document}
